%% LyX 2.4.4 created this file.  For more info, see https://www.lyx.org/.
%% Do not edit unless you really know what you are doing.
\documentclass[english]{beamer}
\usepackage[T1]{fontenc}
\usepackage[latin9]{inputenc}
\setcounter{secnumdepth}{3}
\setcounter{tocdepth}{3}
\usepackage{units}
\usepackage{amssymb}

\makeatletter
%%%%%%%%%%%%%%%%%%%%%%%%%%%%%% Textclass specific LaTeX commands.
% this default might be overridden by plain title style
\newcommand\makebeamertitle{\frame{\maketitle}}%
% (ERT) argument for the TOC
\AtBeginDocument{%
  \let\origtableofcontents=\tableofcontents
  \def\tableofcontents{\@ifnextchar[{\origtableofcontents}{\gobbletableofcontents}}
  \def\gobbletableofcontents#1{\origtableofcontents}
}

%%%%%%%%%%%%%%%%%%%%%%%%%%%%%% User specified LaTeX commands.
\beamertemplateshadingbackground{red!5}{structure!5}

\usepackage{beamerthemeshadow}
\usepackage{pgfnodes,pgfarrows,pgfheaps}
\usepackage{caption}

\beamertemplatetransparentcovereddynamicmedium

\pgfdeclareimage[width=0.6cm]{icsi-logo}{beamer-icsi-logo}
\logo{\pgfuseimage{icsi-logo}}

\newcommand{\Class}[1]{\operatorname{\mathchoice
  {\text{\small #1}}
  {\text{\small #1}}
  {\text{#1}}
  {\text{#1}}}}

\newcommand{\Lang}[1]{\operatorname{\text{\textsc{#1}}}}

% This gets defined by beamerbasecolor.sty, but only at the beginning of
% the document
\colorlet{averagebackgroundcolor}{normal text.bg}

\newcommand{\tape}[3]{%
  \color{structure!30!averagebackgroundcolor}
  \pgfmoveto{\pgfxy(-0.5,0)}
  \pgflineto{\pgfxy(-0.6,0.1)}
  \pgflineto{\pgfxy(-0.4,0.2)}
  \pgflineto{\pgfxy(-0.6,0.3)}
  \pgflineto{\pgfxy(-0.4,0.4)}
  \pgflineto{\pgfxy(-0.5,0.5)}
  \pgflineto{\pgfxy(4,0.5)}
  \pgflineto{\pgfxy(4.1,0.4)}
  \pgflineto{\pgfxy(3.9,0.3)}
  \pgflineto{\pgfxy(4.1,0.2)}
  \pgflineto{\pgfxy(3.9,0.1)}
  \pgflineto{\pgfxy(4,0)}
  \pgfclosepath
  \pgffill

  \color{structure}  
  \pgfputat{\pgfxy(0,0.7)}{\pgfbox[left,base]{#1}}
  \pgfputat{\pgfxy(0,-0.1)}{\pgfbox[left,top]{#2}}

  \color{black}
  \pgfputat{\pgfxy(-.1,0.25)}{\pgfbox[left,center]{\texttt{#3}}}%
}

\newcommand{\shorttape}[3]{%
  \color{structure!30!averagebackgroundcolor}
  \pgfmoveto{\pgfxy(-0.5,0)}
  \pgflineto{\pgfxy(-0.6,0.1)}
  \pgflineto{\pgfxy(-0.4,0.2)}
  \pgflineto{\pgfxy(-0.6,0.3)}
  \pgflineto{\pgfxy(-0.4,0.4)}
  \pgflineto{\pgfxy(-0.5,0.5)}
  \pgflineto{\pgfxy(1,0.5)}
  \pgflineto{\pgfxy(1.1,0.4)}
  \pgflineto{\pgfxy(0.9,0.3)}
  \pgflineto{\pgfxy(1.1,0.2)}
  \pgflineto{\pgfxy(0.9,0.1)}
  \pgflineto{\pgfxy(1,0)}
  \pgfclosepath
  \pgffill

  \color{structure}  
  \pgfputat{\pgfxy(0.25,0.7)}{\pgfbox[center,base]{#1}}
  \pgfputat{\pgfxy(0.25,-0.1)}{\pgfbox[center,top]{#2}}

  \color{black}
  \pgfputat{\pgfxy(-.1,0.25)}{\pgfbox[left,center]{\texttt{#3}}}%
}

\pgfdeclareverticalshading{heap1}{\the\paperwidth}%
  {color(0pt)=(black); color(1cm)=(structure!65!white)}
\pgfdeclareverticalshading{heap2}{\the\paperwidth}%
  {color(0pt)=(black); color(1cm)=(structure!55!white)}
\pgfdeclareverticalshading{heap3}{\the\paperwidth}%
  {color(0pt)=(black); color(1cm)=(structure!45!white)}
\pgfdeclareverticalshading{heap4}{\the\paperwidth}%
  {color(0pt)=(black); color(1cm)=(structure!35!white)}
\pgfdeclareverticalshading{heap5}{\the\paperwidth}%
  {color(0pt)=(black); color(1cm)=(structure!25!white)}
\pgfdeclareverticalshading{heap6}{\the\paperwidth}%
  {color(0pt)=(black); color(1cm)=(red!35!white)}

\newcommand{\heap}[5]{%
  \begin{pgfscope}
    \color{#4}
    \pgfheappath{\pgfxy(0,#1)}{\pgfxy(-#2,0)}{\pgfxy(#2,0)}
    \pgfclip
    \begin{pgfmagnify}{1}{#1}
      \pgfputat{\pgfpoint{-.5\paperwidth}{0pt}}{\pgfbox[left,base]{\pgfuseshading{heap#5}}}
    \end{pgfmagnify}
  \end{pgfscope}
  %\pgffill
  
  \color{#4}
  \pgfheappath{\pgfxy(0,#1)}{\pgfxy(-#2,0)}{\pgfxy(#2,0)}
  \pgfstroke

  \color{white}
  \pgfheaplabel{\pgfxy(0,#1)}{#3}%
}


\pgfdeclareradialshading{graphnode}
  {\pgfpoint{-3pt}{3.6pt}}%
  {color(0cm)=(beamerexample!15);
    color(2.63pt)=(beamerexample!75);
    color(5.26pt)=(beamerexample!70!black);
    color(7.6pt)=(beamerexample!50!black);
    color(8pt)=(beamerexample!10!averagebackgroundcolor)}

\newcommand{\graphnode}[2]{
  \pgfnodecircle{#1}[virtual]{#2}{8pt}
  \pgfputat{#2}{\pgfbox[center,center]{\pgfuseshading{graphnode}}}
}

%\setbeamerfont{author}{size=\Huge} 
%\setbeamerfont{date}{size=\small} 
%\setbeamerfont{exampletext}{series=\bfseries}
%\setbeamerfont{frametitle}{series=\bfseries,size=\large}
%\setbeamerfont{framesubtitle}{series=\mdseries,size=\large}
%\setbeamerfont{section}{series=\bfseries} 
\setbeamerfont{title}{series=\bfseries}

\setbeamercovered{invisible}

\makeatother

\usepackage{babel}
\begin{document}
\title{Calc III: 9.2 Sequences}
\date{{}}

\makebeamertitle
\author{}

\institute{}

\subtitle{}
\begin{frame}{Limit of a Sequence}

\textbf{9.1 55) ${\displaystyle \sum_{k=1}^{n}=\frac{1}{4k^{2}-1}=\frac{n}{2n+1}}$}.
An \emph{explicit formula}\textbf{ }for a sequence of partial sums
has the advantage of being a function. 

\uncover<2->{${\displaystyle \sum_{k=1}^{\infty}\frac{1}{4k^{2}-1}=}{\displaystyle \lim_{n\rightarrow\infty}\frac{n}{2n+1}=\frac{1}{2}}$}
\begin{block}<3->{Recall}

\begin{itemize}
\item <4->\textbf{l'H\^{o}pital's}: ${\displaystyle \lim_{n\rightarrow c}\frac{f(x)}{g(x)}={\displaystyle \lim_{n\rightarrow c}\frac{f'(x)}{g'(x)}}}$
if ${\displaystyle \lim_{n\rightarrow\infty}\frac{f'(x)}{g'(x)}}$
exists and $g'(x)\neq0\,\forall x\in I\setminus c$. 
\item <5->$f(x)^{g(x)}=e^{g(x)\ln f(x)}\Rightarrow{\displaystyle \lim_{n\rightarrow a}f(x)^{g(x)}=e^{{\displaystyle (\lim_{n\rightarrow a}g(x)\ln f(x))}}}$
\item <6->\textbf{Chain rule:} $(f\circ g)'(x)=f'(g(x))g'(x)$
\end{itemize}
\end{block}
\end{frame}

\begin{frame}{}

\textbf{9.2 16) }${\displaystyle \lim_{n\rightarrow\infty}\left(\frac{n}{n+5}\right)^{n}}$ 
\begin{itemize}
\item <2->[]${\color{blue}{\displaystyle \lim_{n\rightarrow\infty}}\left(\frac{n}{n+5}\right)^{n}=e^{\lim n\ln(\frac{n}{n+5})}}\mathrel{\color{blue}\Rightarrow}$
we need to find ${\displaystyle \lim_{n\rightarrow\infty}n\ln\left(\frac{n}{n+5}\right)}={\displaystyle \lim_{n\rightarrow\infty}}\frac{\ln\left(\frac{n}{n+5}\right)}{\nicefrac{1}{n}}$
\item <3->[]\textcolor{blue}{By l'H\^{o}pital's we need to find,} ${\displaystyle \lim_{n\rightarrow\infty}\frac{\ln\left(\frac{n}{n+5}\right)}{\nicefrac{1}{n}}=\lim_{n\rightarrow\infty}\frac{\left(\frac{n+5}{n}\right)\left(\frac{5}{(n+5)^{2}}\right)}{-\nicefrac{1}{n^{2}}}=-5}$. 
\item <4->[]\textcolor{blue}{By $(f\circ g)'(x)=f'(g(x))g'(x)$ and the
quotient rule.}
\item <5->[]So then \textrm{\uncover<6->{$\lim_{n\rightarrow\infty}\left(\frac{n}{n+5}\right)^{n}=e^{-5}.$}}
\end{itemize}
\end{frame}

\begin{frame}{}

\begin{block}{Squeeze Theorem for sequences}

Let $\{a_{n}\},\{b_{n}\},$ and $\{c_{n}\}$ are sequences such that
$a_{n}\leq b_{n}\leq c_{n}$ for all $n>N$ for some number $N$.
If ${\displaystyle \lim_{n\rightarrow\infty}a_{n}}={\displaystyle \lim_{n\rightarrow\infty}c_{n}=L}$
then ${\displaystyle \lim_{n\rightarrow\infty}b_{n}=L}$.
\end{block}
\begin{enumerate}
\item <2->[]\textbf{9.2 45)}\uncover<3->{$\{\frac{-2\pi}{n^{3}+4}\}<$}$\{\frac{2\tan^{-1}n}{n^{3}+4}\}$\uncover<3->{$<\{\frac{2\pi}{n^{3}+4}\}$}
\item <4->[]${\displaystyle \lim_{n\rightarrow\infty}\frac{-2\pi}{n^{3}+4}=\lim_{n\rightarrow\infty}\frac{2\pi}{n^{3}+4}=0}$
\item <5->[]So then ${\displaystyle \lim_{n\rightarrow\infty}\frac{2\tan^{-1}n}{n^{3}+4}=0}.$
\end{enumerate}
\end{frame}

\begin{frame}{}

\begin{block}{}

Let $\{b_{n}\}$ grows faster $\{a_{n}\}$ (notated $\{a_{n}\}\ll\{b_{n}\}$).
Then ${\displaystyle \lim_{n\rightarrow\infty}\frac{\{a_{n}\}}{\{b_{n}\}}=0}$
and ${\displaystyle \lim_{n\rightarrow\infty}\frac{\{b_{n}\}}{\{a_{n}\}}=\infty}$. 
\end{block}
\uncover<2->{\textbf{ex ) ${\displaystyle \lim_{n\rightarrow\infty}\frac{x^{2}}{x^{3}}=0}$}.}
\begin{block}<3->{Growth rate of sequences}

$\{\ln n\}\ll\{n^{p}\}\ll\{n^{p}\ln n\}\ll\{n^{p+s}\}\ll\{b^{n}\}\ll\{n!\}\ll\{n^{n}\}$
for $p,s\in\mathbb{R}$ and $b>1$.
\end{block}
\begin{enumerate}
\item <4->[]\textbf{ex) }${\displaystyle \lim_{n\rightarrow\infty}\frac{n^{\sqrt{7}}}{n^{\sqrt{7}}\ln n}}$\uncover<5->{$=0$} 
\end{enumerate}
\end{frame}

\begin{frame}{Limits}

\begin{block}{Limit of a Sequence}

$\{a_{n}\}$ converges to $L$ if given any $\epsilon>0$ there exists
an $N$ such that $|a_{n}-L|<\epsilon$ whenever $n>N$.
\end{block}
\begin{enumerate}
\item <2->[]Prove (using the definition) that ${\displaystyle \lim_{n\rightarrow\infty}\frac{1}{n}=0}$.
\item <3->[](Want to choose an $N$ s.t. $|\frac{1}{n}-0|<\frac{1}{N}<\epsilon$.)
\item <4->[]Proof: Let $N>\frac{1}{\epsilon}$. So then if $n>N$,
\item <5->[]$|\frac{1}{n}-0|=\frac{1}{n}<\frac{1}{N}$\uncover<6->{$<\epsilon\,$}\uncover<7->{$\square$.} 
\end{enumerate}
\end{frame}

\begin{frame}{}

\textbf{HW 9.2: }4, 9, 18, 20, 29, 46, 48, 51-56, 58 
\end{frame}

\end{document}
